\begin{abstract}
Waste sorting needs detectors that find every item in a photograph, yet much published work classifies single objects on studio backgrounds, and metaheuristic tuning claims rarely include a random-search baseline or confidence intervals. We ask whether particle swarm optimization (PSO) trains a better lightweight waste detector than random search or library defaults at an equal budget. Six training hyperparameters of YOLOv8n were set by the defaults, the best of 32 random samples, or 8 particles over 4 rounds, each search spending 32 short proxy trainings; a check fixed in advance confirmed the proxy ranks settings like longer training (Spearman correlation 0.976, first generation). Each chosen setting was trained and tested once on held-out data, with paired-bootstrap 95\% confidence intervals. The first generation used studio images; the second added real-world photographs and a test-only Indian street set. On studio data, defaults, random search and PSO tie at 0.676, 0.672 and 0.668 mean average precision at intersection-over-union 0.5 (\mapfifty). With real-world training data, both searches significantly beat the defaults on the mixed test set (by 0.017 and 0.020 \mapfifty) but not each other; on real-world photographs (0.298--0.310) no difference is significant, and Indian photographs score only 0.023--0.027. On the same test images, real-world training data raised real-world \mapfifty from 0.038--0.096 to 0.338--0.353 and cost 0.074--0.105 on studio images, all significantly. The optimizer matters little and the data matter most; tuning claims need equal-budget random-search baselines and confidence intervals.
\end{abstract}

\begin{IEEEkeywords}
Waste detection, particle swarm optimization, hyperparameter optimization, YOLOv8, random search, real-world generalization, India.
\end{IEEEkeywords}
